%
\subsection*{Desplazamiento, equilibrio y energía}
\label{CH05-B01}
Fijado el sentido positivo, $m x''$ es la suma de fuerzas. Medido desde
el equilibrio estático, el oscilador lineal habitual cumple
\[
 m x''+c x'+kx=F(t),\qquad m>0,\quad k>0,\quad c\ge0.
\]
Aquí $k$ se mide en fuerza/longitud y $c$ en fuerza/velocidad. Si el
equilibrio ya incorpora el peso, no se añade éste de nuevo. Un evento
inicial exige precisar los datos inmediatamente después de él.

La frecuencia natural angular es $\omega_0=\sqrt{k/m}$, en radianes por
unidad de tiempo; la frecuencia ordinaria es $f_0=\omega_0/(2\pi)$, en
hercios si el tiempo se mide en segundos. La energía
$E=mx'^2/2+kx^2/2$ satisface $E'=F x'-c x'^2$.
Esta identidad comprueba signos y disipación. Recuérdense la segunda ley
de Newton y la derivación de productos. Consulta: solución de
\exref{M2-E0017}.

%
\subsection*{Regímenes, amplitud y envolvente}
\label{CH05-B02}
Las raíces de $mr^2+cr+k=0$ clasifican la respuesta libre. Con $c>0$,
$c^2<4mk$ da oscilación amortiguada; $c^2=4mk$, amortiguamiento crítico;
$c^2>4mk$, sobreamortiguamiento. En el primer caso,
\[
 x(t)=R e^{-ct/(2m)}\cos(\omega_d t-\phi),\qquad
 \omega_d=\sqrt{k/m-c^2/(4m^2)}.
\]
Las curvas $\pm R e^{-ct/(2m)}$ acotan la respuesta; sus valores no son
automáticamente máximos de la trayectoria, que se localizan con $x'=0$.
Tampoco una envolvente de una familia es necesariamente una trayectoria
de esa familia.

En un RLC serie ideal, la carga $q$ y la corriente $i=q'$ cumplen
$Lq''+Rq'+q/C=V(t)$; las relaciones de inductor, resistencia y
condensador fijan los signos según el circuito. Recuérdense raíces
complejas y trigonometría. Consulta: solución de \exref{M2-E0374}.

%
\subsection*{Forzamiento periódico y resonancia}
\label{CH05-B03}
Para $F(t)=F_0\cos\omega t$, $F_0>0$, una respuesta periódica puede
escribirse $A\cos(\omega t-\delta)$, donde
\[
 A=\frac{F_0}{\sqrt{(k-m\omega^2)^2+(c\omega)^2}},\qquad
 \cos\delta=\frac{k-m\omega^2}{D},\quad
 \sin\delta=\frac{c\omega}{D},
\]
con $D$ igual al denominador y $D\ne0$. Ambas razones fijan el cuadrante
del desfase. Con $c>0$, el transitorio decae. Con $c=0$ y
$\omega=\omega_0$, se necesita una particular resonante, proporcional
a $t\sin(\omega_0t)$.

El PVI conserva sus datos: al variar la frecuencia, el límite resonante
se toma sobre la respuesta completa, incluida la homogénea. En un apoyo
móvil $b(t)$, la elongación y la velocidad relativa son $x-b$ y $x'-b'$;
el balance determina dónde aparecen. Consulta: solución de
\exref{M2-E0117}.

%
\subsection*{Identificación y garantías}
\label{CH05-B04}
Una inversión algebraica puede producir varias ramas. Manténganse todas
las que respetan unidades, signos y dominio físico hasta aplicar los datos
que permiten distinguirlas. Ni un dato aislado ni un régimen final
identifican necesariamente todos los parámetros o la ley resistente.

Un intervalo garantizado describe valores posibles; no asigna
probabilidades. Para afirmar una propiedad para todos los parámetros
admisibles hace falta una cota válida en todo el conjunto. Un solo caso
admisible que la incumpla refuta esa garantía. Al buscar máximos, indíquese
si se optimiza sobre el tiempo de una trayectoria, sobre parámetros a
tiempo fijo o sobre ambos: las preguntas son diferentes. Recuérdense
desigualdades y restricciones físicas. Consultas: soluciones de
\exref{M2-E0038} y \exref{M2-E0374}.
